% `tex.web` §873: what the third pass actually *does* with `\emergencystretch`.
%
%     else begin if tracing_paragraphs>0 then print_nl("@emergencypass");
%       background[2]:=background[2]+emergency_stretch; final_pass:=true;
%       end;
%
% `background[2]` is the normal-order stretch every line starts with —
% `\leftskip` and `\rightskip` — and the emergency pass adds `\emergencystretch`
% to it. So the pass does not look for different breaks; it makes the *same*
% breaks less bad, which is how a paragraph that would otherwise go overfull
% comes out merely loose.
%
% We ran the third pass with the parameter deciding only *whether* it happened
% and never widening anything, so it scored exactly what the second pass had
% and fell through to §851's artificial demerits — `b=*` in the narration, and
% a forced break. TeX gets `b=99` and breaks a word earlier.
%
% LaTeX sets `\emergencystretch` in `\sloppy` and in every `\raggedbottom`
% class, so this is not an exotic path.
\catcode`\{=1 \catcode`\}=2
\tracingonline=1 \tracingparagraphs=1 \showboxbreadth=99 \showboxdepth=99
\hsize=90pt \parindent=0pt \tolerance=100 \pretolerance=-1
\font\r=cmr10 \r
% The box warnings are not what this case is about, and the two passes differ
% in which of them they raise.
\hbadness=10000 \vbadness=10000
\showthe\escapechar

% No emergency stretch: the second pass is the final one, and a word that will
% not fit forces a break.
\setbox0\vbox{AAAAAAAA BBBBBBBB CCCCCCCCCCCCCCCCCC DD\par}
\showbox0

% ...and with it, a third pass whose lines have 30pt more to give.
\emergencystretch=30pt
\setbox1\vbox{AAAAAAAA BBBBBBBB CCCCCCCCCCCCCCCCCC DD\par}
\showbox1

% Back to zero: no third pass, and no `@emergencypass` line.
\emergencystretch=0pt
\setbox2\vbox{AAAAAAAA BBBBBBBB CCCCCCCCCCCCCCCCCC DD\par}
\showbox2
\showthe\escapechar
\end
